%% ============================================================
%% SAMBP TSA Meeting — Speaker Notes (Standalone Document)
%% Author  : Anoop V Eluvathingal, IIT Madras
%% Compile : pdflatex sambp_speaker_notes.tex
%% ============================================================
\documentclass[12pt, a4paper]{article}
\usepackage[T1]{fontenc}
\usepackage[utf8]{inputenc}
\usepackage{mathptmx}
\usepackage[a4paper, top=20mm, bottom=20mm, left=22mm, right=22mm]{geometry}
\usepackage{xcolor}
\usepackage{titlesec}
\usepackage{enumitem}
\usepackage{booktabs}
\usepackage{microtype}
\usepackage{setspace}
\usepackage{fancyhdr}
\usepackage{amsmath}
\usepackage{fontawesome5}
\usepackage[colorlinks=true, linkcolor=iitmMaroon, urlcolor=iitmMaroon]{hyperref}

\definecolor{iitmMaroon}{RGB}{128,0,0}
\definecolor{iitmGold}  {RGB}{179,145,25}
\definecolor{noteBg}    {RGB}{252,250,245}
\definecolor{slideNum}  {RGB}{0,60,130}

\setstretch{1.3}
\setlength{\parindent}{0pt}
\setlength{\parskip}{5pt}

%% Section formatting — slide title bar
\newcommand{\secbox}[1]{%
  \colorbox{iitmMaroon}{\parbox{\dimexpr\textwidth-2\fboxsep\relax}{%
    \strut\quad\color{white}\large\bfseries\thesection\enspace|\enspace #1\strut}}}
\titleformat{\section}[block]{\normalfont}{}{0pt}{\secbox}[\vspace{2pt}]

\titlespacing{\section}{0pt}{14pt}{6pt}

\pagestyle{fancy}
\fancyhf{}
\renewcommand{\headrulewidth}{0.4pt}
\fancyhead[L]{\small\color{iitmMaroon}\textit{SAMBP TSA Meeting --- Speaker Notes}}
\fancyhead[R]{\small\color{iitmMaroon}\textit{Anoop V Eluvathingal, IIT Madras}}
\fancyfoot[C]{\small\thepage}

%% Slide note box
\usepackage{tcolorbox}
\tcbuselibrary{skins}
\newenvironment{notebox}{%
  \begin{tcolorbox}[
    colback=noteBg,
    colframe=iitmMaroon!40,
    boxrule=0.6pt,
    arc=3pt,
    left=8pt, right=8pt, top=6pt, bottom=6pt
  ]
  \setstretch{1.4}\small
}{%
  \end{tcolorbox}
}

\newcommand{\slidetag}[2]{%
  \noindent
  \textcolor{slideNum}{\footnotesize\bfseries Slide #1} \hfill
  \textcolor{iitmGold!80!black}{\footnotesize\itshape Approx.\ #2}\\[2pt]
}

\begin{document}

%% ── Title page ──────────────────────────────────────────────
\begin{center}
  {\color{iitmMaroon}\rule{\linewidth}{2pt}}\\[0.5cm]
  {\LARGE\bfseries\color{iitmMaroon}
    Advanced Protection for Distributed\\[0.2cm]
    Power Systems and Microgrids}\\[0.4cm]
  {\large Thesis Submission Approval (TSA) Meeting}\\[0.3cm]
  {\large\bfseries Speaker Notes}\\[0.5cm]
  {\normalsize Anoop V.\ Eluvathingal \quad|\quad
               Supervisor: Prof.\ K.\ Shanti Swarup}\\[0.1cm]
  {\normalsize Dept.\ of Electrical Engineering, IIT Madras \quad|\quad June 2026}\\[0.5cm]
  {\color{iitmMaroon}\rule{\linewidth}{2pt}}
\end{center}

\vspace{0.3cm}
\begin{tcolorbox}[colback=iitmMaroon!8, colframe=iitmMaroon, arc=3pt,
                  boxrule=0.8pt, left=8pt, right=8pt, top=6pt, bottom=6pt]
  \small
  \textbf{How to use these notes:}
  \begin{itemize}[leftmargin=1.5em, itemsep=2pt, topsep=3pt]
    \item Each section corresponds to one slide. Total: 12 numbered + title + thank-you.
    \item Estimated timing is shown in the top-right of each slide section.
    \item Key phrases to emphasise are shown in \textbf{bold}.
    \item Appendix slides (A--I) are backup --- use only during Q\&A.
    \item Total target delivery time: \textbf{18--22 minutes} for the 12 content slides.
  \end{itemize}
\end{tcolorbox}

\newpage
\tableofcontents
\newpage

%% ============================================================
\section{Title Slide}

\slidetag{---}{1 min}
\begin{notebox}
Good morning / good afternoon, respected committee members.
I am \textbf{Anoop V.\ Eluvathingal}, a PhD candidate in the Department of Electrical
Engineering at IIT Madras, working under the supervision of
\textbf{Prof.\ K.\ Shanti Swarup}.

Today I present my thesis titled \textbf{``Advanced Protection for Distributed Power
Systems and Microgrids''} for the Thesis Submission Approval.

My work addresses a fundamental challenge in modern power grids: how do we protect
networks that are increasingly dominated by \textbf{inverter-based resources} (IBRs),
which behave fundamentally differently from the synchronous generators our relay
algorithms were designed for?

Over the next 20 minutes I will walk you through the research background, the six key
contributions, and our validated results. Backup slides covering each thesis chapter
in detail are available after the last numbered slide for the Q\&A.
\end{notebox}

%% ============================================================
\section{Outline}

\slidetag{1 of 12}{0.5 min}
\begin{notebox}
I have organised this presentation into \textbf{12 sections}.

We begin with the background and motivation, then critically examine what the state of
the art can and cannot do. The core of the talk covers six key contributions spread
across three slides --- one per two contributions --- followed by the thesis organisation,
publications, and a closing summary with recommended next steps.

Please hold questions until the end; I am happy to go to any backup chapter slide
during discussion.
\end{notebox}

%% ============================================================
\section{Background and Context}

\slidetag{2 of 12}{1.5 min}
\begin{notebox}
The left chart tells the story of the generation mix transformation.
By 2030, IBRs are projected to supply \textbf{72\% of generation} while synchronous
machines drop to 28\%. This is already happening in India's renewable-rich states.

The right chart shows the \textbf{critical technical consequence}.
A synchronous generator produces 8--10 per-unit fault current.
An inverter-based resource is \textbf{software-limited to 1.1--2.0 pu}.

The red dotted line is a typical overcurrent relay pickup set for synchronous machines.
When an IBR feeds a fault, the current simply does not reach the pickup threshold ---
the relay does not see a fault at all.

This is the physical root of every failure mode I will show next.
\textbf{IEEE C37.300-2024}, the new CPC systems standard, explicitly mandates an
adaptive supervisory layer for exactly this reason.
\end{notebox}

%% ============================================================
\section{Research Motivation}

\slidetag{3 of 12}{2 min}
\begin{notebox}
The 12 failure modes on the right are not theoretical --- they are \textbf{systematically
documented from grid incident reports and protection literature}.

The diagram on the left illustrates the most insidious mode: \textbf{sympathetic tripping}.
A fault on Feeder 1 causes the IBR on Feeder 2 to push current into the fault through
the common bus. R2 on the \emph{healthy} feeder sees this elevated current and trips
unnecessarily, disconnecting healthy load.

The table spans all major protection elements: overcurrent 51, distance 21, directional
67, line differential 87L, islanding, pilot schemes, and CUSUM-based monitoring.

\textbf{All 12 modes trace to one root cause}: relay parameters are static while IBR
behaviour is dynamic. The conclusion is clear --- we need an adaptive supervisory layer
that continuously refreshes relay settings from real-time estimates.
\end{notebox}

%% ============================================================
\section{State-of-the-Art \& Limitations}

\slidetag{4 of 12}{2 min}
\begin{notebox}
The sequence network diagram reveals a critical gap in the literature.
In the positive-sequence network, IBR current is limited to $k_\text{ibr}$.
But in the \textbf{negative-sequence network}, the IBR presents zero impedance ---
it does not limit negative-sequence current in the same way.

This leads to \textbf{Proposition 4.2} of the thesis:
\emph{when IBR suppresses $I_2$, the 87LN operate ratio $\rho_2 = 2$ always,
regardless of $k_\text{ibr}$}. IBR suppression is therefore a \textbf{protection asset}
for negative-sequence elements --- an insight not previously reported.

The gap table summarises seven protection areas. Each row represents a specific chapter
of this thesis. The final row highlights the absence of any closed-loop, multi-tier
coordination framework --- which is exactly what SAMBP provides.
\end{notebox}

%% ============================================================
\section{Research Objectives}

\slidetag{5 of 12}{1.5 min}
\begin{notebox}
The six objectives map \textbf{directly to the six contributions}.
O1--O3 address local relay adaptation: line differential, distance, and microgrid
islanding. O4 and O5 address the monitoring layer: real-time parameter estimation and
physics-informed machine learning. O6 closes the loop at the coordination layer.

The study network on the right is the common testbed for all six contributions: 10 buses,
7 IBR units, and 3 PMUs, configured to represent a realistic distribution-level active
network. All Monte Carlo campaigns were run on this network.

The $k_\text{ibr}$ parameter --- defined as the IBR current limit normalised by rated
current --- ranges from 0.03 to 0.85 pu across the study scenarios, covering everything
from near-zero to near-rated IBR output.
\end{notebox}

%% ============================================================
\section{Key Contributions (1 of 3)}

\slidetag{6 of 12}{2 min}
\begin{notebox}
\textbf{C1 --- Adaptive 87L:}
The core result is the adaptive slope law $k_m^*(k_\text{ibr})$.
A fixed slope fails at both extremes of $k_\text{ibr}$: too sensitive at high
$k_\text{ibr}$ causing false trips; too insensitive at low $k_\text{ibr}$ causing missed
trips. Our law tracks the through-fault operate ratio in real time.
We also introduce \textbf{87LN} for zero-sequence sensitivity in pure-IBR networks,
and \textbf{TDCS} --- Transient Differential Catch Scheme --- as a backstop for the
sub-threshold blind spot.

The bar chart confirms \textbf{100\% detection for any GFL/GFM mix}. The 67Q element
activates only above 50\% GFM fraction when negative-sequence directional discrimination
becomes meaningful.

\textbf{C2 --- Adaptive distance:}
IBR phase-angle injection shifts the apparent impedance outside the relay mho circle.
The adaptive quadrilateral reach corrects for this with a \textbf{$+7.6$ percentage-point
detection gain} validated over 10{,}000 Monte Carlo trials.
\end{notebox}

%% ============================================================
\section{Key Contributions (2 of 3)}

\slidetag{7 of 12}{2 min}
\begin{notebox}
This diagram is the \textbf{architectural summary} of the entire thesis. It is important
to read this correctly.

\textbf{Tier 1} (blue, left) is the local relay layer --- adaptive 87L, sequence
elements, and IBR modelling. Chapters 2--4.

\textbf{Tier 2} (lighter blue) extends to distance and microgrid protection.
Chapters 5--6.

\textbf{Tiers 3 and 4} (green) form the monitoring layer --- EKF digital twin and PINN
classifier. Chapters 7--8.

\textbf{Tier 5} (maroon) is the Wide-Area Protection Coordinator which sends updated
settings back to the relays. Chapter 9.

\textbf{Tier 6} (gold) is the Digital Twin validation and audit framework spanning
Chapters 7--9.

The arrows between tiers show the upward information flow. This is a
\textbf{closed feedback loop}, not a hierarchical command system. The relay can always
act on its own physics; the upper tiers only refine it.
\end{notebox}

%% ============================================================
\section{Key Contributions (3 of 3)}

\slidetag{8 of 12}{2 min}
\begin{notebox}
\textbf{C3 --- Islanding detection:}
The combined ROCOF plus $V_2/V_1$ criterion closes the non-detection zone to below
\textbf{1\% of the load range}, while maintaining IEEE 1547-2018 ride-through
compliance. Setting-group switching via GOOSE completes in 50 ms.

\textbf{C4 --- EKF digital twin:}
The dual-core shadow-register design ensures the TRIP core always uses the most recent
committed estimate. The EKF estimation cycle of 20 ms is completely decoupled from the
relay's 50 $\mu$s trip decision.

\textbf{C5 --- PINN:}
Physics constraints (KVL, $I_1 = I_2^*$ symmetry, and current limiting) are embedded
directly in the loss function. At \textbf{96.1\% 5-class accuracy} with zero physics
violations and 3.7 $\mu$s inference, this is directly deployable in real relay firmware.

\textbf{C6 --- WAPC:}
The closed-loop protocol runs CUSUM monitoring, parameter computation, digital twin
validation, MMS upload, and relay activation in a total of \textbf{31 minutes}.
HMAC-SHA256 authentication reduces GOOSE message spoofing risk by \textbf{54.9\%}.

\textit{Draw attention to the validation metrics table on the right --- these numbers
are the key takeaways.}
\end{notebox}

%% ============================================================
\section{Thesis Organization}

\slidetag{9 of 12}{1.5 min}
\begin{notebox}
The thesis follows a deliberate progression.
\textbf{Chapters 1--3} establish the foundation: the SAMBP architecture, the physics
of IBR fault currents, and a systematic characterisation of the 12 protection failure
modes.

\textbf{Chapters 4--6} deliver the three local relay contributions --- one chapter per
contribution, each with a full derivation, algorithm, and Monte Carlo validation campaign.

\textbf{Chapters 7--9} address the monitoring and coordination layers. Chapter 9 is the
integrating chapter that ties all contributions together through the Wide-Area Protection
Coordinator and the 31-minute closed-loop protocol.

The 55 SAMBP technical reports are referenced throughout the thesis as the primary
numerical evidence base. They are available for inspection and form a \textbf{complete
audit trail} from algorithm to validated result.
\end{notebox}

%% ============================================================
\section{List of Publications (1 of 2)}

\slidetag{10 of 12}{1.5 min}
\begin{notebox}
The thesis is anchored by a published journal paper in \textit{IEEE Access} (2024).
This paper introduces the concept of adaptive relay logic for resolving the conflict
between line protection and LVRT requirements in medium-voltage solar PV
interconnections --- directly motivating Contribution C1.

The Springer book chapter from 2020 established early foundational concepts in
protection for inverter-interfaced networks.

On the pipeline side, the \textbf{adaptive 87L paper (P1) is currently under review at
\textit{IEEE Transactions on Power Delivery}}. The EKF digital twin paper (P2) is in
preparation for \textit{IEEE Transactions on Power Electronics}. The PINN paper (P3)
targets \textit{IEEE Transactions on Smart Grid}.
\end{notebox}

%% ============================================================
\section{List of Publications (2 of 2)}

\slidetag{11 of 12}{1 min}
\begin{notebox}
Papers P4 and P5 are planned. P4 on the Wide-Area Protection Coordinator targets
\textit{IEEE Transactions on Power Systems}. P5 on the HIL validation is planned for
the IEEE PES General Meeting.

The four conference papers from 2017--2018 represent the investigative phase of the
research that identified the key protection gaps this thesis resolves. They also
demonstrate the evolution from symmetrical-component islanding relays to the full
SAMBP framework.

In total: \textbf{1 published journal paper}, 1 book chapter, 4 conference papers, and
\textbf{5 journal papers in the submission pipeline}. The 55 SAMBP technical reports
provide the full derivation and simulation evidence base for all claims.
\end{notebox}

%% ============================================================
\section{Summary \& Next Steps}

\slidetag{12 of 12}{2 min}
\begin{notebox}
To summarise: this thesis delivers a \textbf{complete six-tier self-adaptive protection
framework} for IBR-dominated power networks.

The six contributions collectively \textbf{eliminate all 12 failure modes} identified in
the motivation. The validation is rigorous: 5{,}000 Monte Carlo trials per group, 100
hardware-in-loop scenarios, and physics consistency verified for every PINN output.

Three open research gaps define the future roadmap. The most technically critical is
\textbf{Gap G4} --- GFM virtual-inertia transient modelling beyond 150 ms, which limits
the applicability of the Stage-1 Park model.

My most immediate practical recommendation: \textbf{begin deploying adaptive 87L
protection before IBR penetration in any feeder zone exceeds 40\%}.

I believe SAMBP provides the \textbf{first complete, closed-loop, self-adaptive
protection architecture} for the next generation of power networks. I am grateful for
the guidance of Prof.\ K.\ Shanti Swarup and the committee throughout this journey.

\textit{Pause here for 2--3 seconds. Then: ``Thank you.''}
\end{notebox}

%% ============================================================
\section*{Thank You / Q\&A}
\addcontentsline{toc}{section}{Thank You / Q\&A}

\slidetag{---}{open}
\begin{notebox}
I am happy to address any questions.

Backup slides are available for each chapter (\textbf{Appendix A through I}) covering:
SAMBP architecture, IBR modelling, protection failure modes, adaptive 87L detail,
distance relay, microgrid islanding, digital twin, PINN classifier, and WAPC.

I can also navigate directly to any of the 55 SAMBP technical reports if the committee
wishes to examine specific derivations or simulation data.

\textit{Common questions and suggested responses:}
\begin{itemize}[leftmargin=1.5em, itemsep=4pt]
  \item \textbf{``How does TDCS avoid false alarms during energisation?''}\\
        The TDCS threshold $\Delta_{thr} = 0.04$ pu is set 15.5\% above the 95th
        percentile CT mismatch envelope. During energisation, current is unidirectional
        so $\Delta I_{rms}$ stays below threshold.

  \item \textbf{``Is the EKF stable under IBR inverter saturation?''}\\
        The EKF uses a linearised model valid for $k_{ibr} \in [0.03, 0.85]$. Beyond
        0.85 pu, the inverse estimator flags low confidence and the gate withholds
        adaptation. The relay then falls back to nominal settings.

  \item \textbf{``What is the computational cost of the PINN on relay hardware?''}\\
        At 3.7 $\mu$s inference time on a standard laptop CPU, the PINN is well within
        the budget for modern numerical relays with DSP co-processors. The model has
        3 hidden layers of 64 neurons --- a small network by current standards.

  \item \textbf{``Why 31 minutes for the WAPC closed loop? Is that acceptable?''}\\
        The 31-minute cycle is for \textbf{setting update}, not for fault protection.
        The relay continues to protect during the update cycle using the last committed
        settings. The 25-minute MMS phase includes a mandatory human operator approval
        step per IEC 61850-90-1. In fully automated deployments this could be reduced
        to under 5 minutes.

  \item \textbf{``How does the system handle simultaneous faults?''}\\
        Each relay acts independently in real time. The WAPC handles coordination
        sequentially after the first fault is cleared, which is the standard assumption
        for CTI-based coordination.
\end{itemize}
\end{notebox}

%% ============================================================
\section*{Appendix Notes (Backup Slides A--I)}
\addcontentsline{toc}{section}{Appendix Notes (Backup Slides A--I)}

\subsection*{Appendix A --- SAMBP Architecture (Ch.\ 1)}
\begin{notebox}
This figure shows the full three-tier SAMBP architecture. The local relay tier handles
all protection actions within 50 $\mu$s. The monitoring tier runs the EKF digital twin
on a 20 ms cycle. The coordination tier runs the CUSUM-to-GOOSE closed loop on a
minute-scale cycle.

Key design principle: \textbf{physics-first, non-blocking}. The relay never waits for
AI or optimisation output to make a trip decision. Upper tiers only improve settings
\emph{between} fault events.
\end{notebox}

\subsection*{Appendix B --- IBR Modelling (Ch.\ 2)}
\begin{notebox}
The GFL inverter control structure shows the current control loop that enforces the
$k_\text{ibr}$ limit. The dq-axis current references are clipped at $k_\text{ibr}$ times
rated current.

The 10-bus study network: IBRs at buses 2--5 (radial feeder) and buses 7--9 (lateral).
Three PMUs at buses 4 and 7 provide wide-area observability for the WAPC.

$k_\text{ibr}$ sweeps from 0.03 to 0.85 pu in all study campaigns,
covering worst-case low-inertia to near-rated injection scenarios.
\end{notebox}

\subsection*{Appendix C --- Protection Challenges (Ch.\ 3)}
\begin{notebox}
The left figure: the 87L failure plane in $I_\text{diff}$--$I_\text{bias}$ space with
fixed slope. Red crosses = missed trips (IBR below slope); orange squares = false trip
risk (IBR drives $I_\text{diff}$ above slope during through-fault conditions).

Core finding --- the operate ratio formula:
$\rho = 2(1-k_\text{ibr}) / (1+k_\text{ibr})$.
When $k_\text{ibr} \to 1$: $\rho \to 0$ (false trip risk).
When $k_\text{ibr} \to 0$: $\rho \to 2$ (missed trip if slope $> 2$).
A fixed slope cannot satisfy both simultaneously.
\end{notebox}

\subsection*{Appendix D --- Adaptive 87L Detail (Ch.\ 4)}
\begin{notebox}
The adaptive slope curve is derived by solving $\rho(k_\text{ibr}) = k_m^*(k_\text{ibr})$
with appropriate security margins.

The formula has three components: a main quadratic term tracking the through-fault
$\rho$, a small linear correction for the high-$k_\text{ibr}$ region, and a 0.05 security
floor.

Bar chart: TDCS catches bins below $k_\text{ibr} = 0.12$ where 87L is too insensitive
even with adaptive slope. Together they achieve TPR = 1.000 across all bins
($N = 5000$ MC trials, TR-26).
\end{notebox}

\subsection*{Appendix E --- Distance Protection (Ch.\ 5)}
\begin{notebox}
R--X locus: IBR angle injection causes the apparent impedance to shift outside the
zone-1 mho circle (overreach) or deep inside (underreach, weak-infeed).

Correction term: $\Delta X = Z_F k_\text{ibr} (I_r/I_A) \sin\beta$ derived analytically
from the sequence network. The adaptive quadrilateral applies
$Z_1^\text{adapt} = Z_1(1 + k_r \hat{k}_\text{ibr})$ using the EKF estimate.

$+7.6$ pp detection gain at $k_\text{ibr} = 0.5$ is representative of typical operating
points.
\end{notebox}

\subsection*{Appendix F --- Microgrid \& Islanding (Ch.\ 6)}
\begin{notebox}
Mode transition FSM: grid-connected GFL mode $\to$ islanded GFM mode, triggered by the
combined detection criterion.

Three-condition logic: ROCOF captures fast frequency excursions; phase angle deviation
catches slower drift events; $V_2/V_1$ detects unbalanced islanding invisible to ROCOF
alone.

End-to-end: 60 ms detection + 4 ms GOOSE + 50 ms relay = \textbf{114 ms total}.
Meets IEEE 1547-2018 ride-through requirements with NDZ below 1\% of load range.
\end{notebox}

\subsection*{Appendix G --- Digital Twin (Ch.\ 7)}
\begin{notebox}
Non-blocking architecture: TRIP core ($\leq 50\,\mu$s) and EKF core (20 ms cycle) run
in parallel, sharing state only via atomic shadow register write.

Three parameters estimated simultaneously:
$\hat{k}_\text{ibr}$ (RMSE 0.024 pu), $\hat{Z}_\text{line}$ (0.5\%), and
$\hat{\alpha}$ (1.1\%).

TDCS margin: 15.5\% safety margin over the 95th percentile CT mismatch envelope.
This sets the TDCS threshold in the 87L implementation.
\end{notebox}

\subsection*{Appendix H --- PINN Classifier (Ch.\ 8)}
\begin{notebox}
Feature importance chart (TR-47): top 3 features are $\alpha$ (fault angle),
$k_\text{ibr}$ (IBR limit), and $R_\text{arc}$ (arc resistance) --- all align with
physics, validating that the PINN has learned physically meaningful representations.

Physics constraint weights: KVL (10) $>$ symmetry $I_1 = I_2^*$ (5) $>$ current
limiting (2). Hierarchy determined through ablation studies.

3.7 $\mu$s inference fits relay firmware budget. Online re-training via CUSUM-SGD
($\leq 21$ cycles) adapts to fleet-level IBR parameter drift automatically.
\end{notebox}

\subsection*{Appendix I --- Wide-Area Protection Coordinator (Ch.\ 9)}
\begin{notebox}
Four WAPC subsystems: fault localisation via PMU $\Delta V$ triangulation (94.5\%),
CTI engine for IEC standard-inverse coordination
($\Delta_\text{CTI} \geq 0.2$ s), GOOSE security via HMAC-SHA256, and HIL validation.

Closed-loop protocol: 0.5 min CUSUM alarm $\to$ 2 min compute $\to$ 1 min DT
validation $\to$ 25 min MMS upload (includes operator approval per IEC 61850-90-1)
$\to$ 2.5 min activate.

54.9\% GOOSE risk reduction from eliminating unauthenticated messages.
HMAC-SHA256 adds only 0.36 ms overhead --- well within the 4 ms GOOSE budget.
\end{notebox}

\end{document}
